% Generated from locale/tr/content/first-order-logic/models-theories/introduction.tex; do not hand-edit.


\olfileid[tr]{fol}{mat}{int}
\olsection{Giriş}

\begin{explain}
Aksiyomatik yöntemin geliştirilmesi bilim tarihinin önemli bir
başarısıdır ve matematik tarihinde özel bir öneme sahiptir. Bir alanın
aksiyomatik olarak geliştirilmesi birçok sorunun açıklığa
kavuşturulmasını içerir: Alan ne hakkındadır? En temel kavramları
nelerdir? Bunlar nasıl ilişkilidir? Alanın bütün kavramları bu temel
kavramlar cinsinden tanımlanabilir mi? Bu kavramlar hangi yasalara uyar
ve uymak zorundadır?

Aksiyomatik yöntem ile mantık birbirleri için biçilmiş kaftandır.
Biçimsel mantık; aksiyomatik kuramları ifade etmek, kuramın
aksiyomlarından teoremleri kesin olarak belirtilmiş biçimde ispatlamak
ve aksiyomları sağlayan bütün sistemlerin özelliklerini sistematik
olarak incelemek için araçlar sağlar.
\end{explain}

\begin{defn}
Bir !!{sentence}s kümesi~$\Gamma$, $\Gamma\Entails !A$ olduğunda
$!A\in\Gamma$ oluyorsa ve ancak bu durumda \emph{kapalı}dır. Bir
!!{sentence}s kümesi~$\Gamma$'nın \emph{kapanışı}
$\Setabs{!A}{\Gamma\Entails !A}$ kümesidir.

$\Gamma$,~$\Delta$'nın kapanışıysa,~$\Gamma$'nın bir tümceler
kümesi~$\Delta$ tarafından \emph{aksiyomlaştırıldığını} söyleriz.
\end{defn}

\begin{explain}
Aksiyomatik bir kuramı, aksiyomları kümesi~$\Delta$ tarafından
aksiyomlaştırılan tümceler kümesi olarak düşünebiliriz. Başka bir
deyişle, incelemek istediğimiz aksiyomatik bilimin temel kavramları
için mantıksal olmayan simgeler içeren bir birinci dereceden dil ile
bilimin temel yasalarını ifade eden bir !!{sentence}s kümesine sahip
olduğumuzda kuramı, bu dilde aksiyomların mantıksal olarak gerektirdiği
bütün !!{sentence}s ile temsil edilmiş sayabiliriz. Bu, tek bir temel
kavram ve basit aksiyomlar içeren kısmi sıralamalar kuramı gibi basit
örneklerden Newton mekaniği gibi karmaşık kuramlara kadar uzanır.

Aksiyomatik yönteme yönelik bu biçimsel yaklaşımı önemli kılan başlıca
mantıksal olgular şunlardır. $\Gamma$ bir kuram için bir aksiyom
sistemi, yani bir tümceler kümesi olsun.
\begin{enumerate}
\item Bir aksiyom sisteminin amaçlanan bir !!{structure}s sınıfını ne
  zaman yakaladığını kesin olarak söyleyebiliriz. Belirli bir
  !!{structure}s sınıfıyla ilgileniyorsak bu sınıfı bir aksiyom
  sistemi~$\Gamma$ ile başarıyla yakalamamız ancak ve ancak söz konusu
  !!{structure}s tam olarak $\Sat M\Gamma$ olan~$\Struct M$ yapılarıysa
  mümkündür.
\item $\Sat M\Gamma$ olduğu halde~$\Struct M$'nin amaçladığımız
  !!{structure}s arasında bulunmaması yüzünden bu konuda başarısız
  olabiliriz. Bu,~$\Struct M$ içinde doğru olmayan aksiyomlar
  eklememize yol açabilir.
\item En azından~$\Gamma$'nın amaçlanan bütün !!{structure}s içinde
  doğru olmasını başarabilirsek $\Gamma\Entails !A$ olduğunda~$!A$
  tümcesi bütün amaçlanan !!{structure}s içinde doğrudur. Böylece
  tümcelerin bütün amaçlanan !!{structure}s içinde doğru olduğunu,
  yalnızca aksiyomların onları mantıksal olarak gerektirdiğini
  göstererek mantıksal araçlarla (örneğin !!{derivation} yöntemleriyle)
  ortaya koyabiliriz.
\item Bazen aklımızda amaçlanan !!{structure}s yoktur; bunun yerine
  aksiyomların kendisinden başlarız. Belirli yasalara uymasını
  istediğimiz bazı temel kavramlarla başlar ve bu yasaları bir aksiyom
  sisteminde kodlarız. Hemen doğrulamak istediğimiz şeylerden biri,
  aksiyomların birbirleriyle çelişmemesidir. Çelişirlerse bu yasalara
  uyan hiçbir kavram olamaz ve tutarsız bir kuram kurmaya çalışmış
  oluruz. Bunun olmadığını~$\Gamma$'nın bir modelini bularak
  doğrulayabiliriz. Kuramımızın modelleri varsa bunları incelemek ve
  model kurmak için mantıksal yöntemleri kullanabiliriz.
\item Aksiyomların bağımsızlığı da önemli bir sorudur. Aksiyomlardan
  biri gerçekte ötekilerin sonucu ve dolayısıyla gereksiz olabilir.
  $\Gamma$ içindeki bir~$!A$ aksiyomunun gereksiz olduğunu
  $\Gamma\setminus\{!A\}\Entails !A$ ispatlayarak gösterebiliriz.
  Bir aksiyomun gereksiz olmadığını ise
  $(\Gamma\setminus\{!A\})\cup\{\lnot !A\}$ kümesinin sağlanabilir
  olduğunu göstererek ispatlayabiliriz. Örneğin paralellik
  postulatının geometrinin öteki aksiyomlarından bağımsız olduğu bu
  biçimde gösterilmiştir.
\item Başka önemli bir soru, bir kuramdaki kavramların
  tanımlanabilirliğidir: Dilin seçimi, bir kuramın modellerinin nelerden
  oluştuğunu belirler. Fakat kuramın her yönünün modellerinde ayrı ayrı
  temsil edilmesi gerekmez. Örneğin her~$\leq$ sıralaması karşılık
  gelen bir katı sıralama~$<$ belirler; biri verildiğinde ötekini
  tanımlayabiliriz. Dolayısıyla böyle bir sıralamayı içeren bir kuramın
  modelinin karşılık gelen katı sıralamayı \emph{ayrıca} içermesi
  gerekmez. Genel olarak bir bağıntı ne zaman ötekiler cinsinden
  tanımlanabilir? Bir bağıntıyı ötekiler cinsinden tanımlamak ne zaman
  olanaksızdır (ve bu nedenle dilin temel kavramlarına eklenmelidir)?
\end{enumerate}
\end{explain}

